\section{讨论与有效性威胁}
\subsection{新增实证改变了哪些判断}
本轮两处源码问题获得了真实进程失败、精确参照与修复输出，说明反例驱动回归在登记子契约上有可操作价值。
其适用性来自可机械验证的目标，不来自Judge给高分。模板成功输出仍可能没有目标产物，任务空值定义也可与工具默认语义不同。
外部评审则没有证明证据提示提升区分能力；三臂引用定位都有效仍只有约一半平衡准确率，反驳以定位率代理质量准确率。
但不能把355次机械CLI检查与86次语义轨迹评审直接作总体优劣比较：二者对象、参照与任务难度不同。

402余额不足造成结构性中断，未运行的是排序较后的任务，不能假设缺失随机。
实际子集又没有可查产物，产物增强机制尚未接受有效外部检验。
本轮的费用停止是真实实验限制，不能删除后改称203例全量成功。
自然工具联合回放也只是事后探索：修复候选让简单路径48/48，原候选中过度过滤反而减少求解。
因此目前没有证据认为更复杂路由总优于简单修复及合理选路。

\subsection{可信性来自可反驳性及独立测量}
质量判别器对明确规则输出资源级反证，对行为输出独立检查结果，且版本和输入指纹可复核。
相比只给文档平均分，它能解释缺陷在哪里，以及哪些请求仍未覆盖。
对照、模块消融、保留失败与未知、冻结选择后再执行，降低了循环论证和只报告有利案例的风险。
网络参照采用与worker不同的CSV读取和精确算术路线，重放还能核对每个选择及产物。
这些提高的是本设计的内部可核验性，不等于独立领域专家已认可语义金标准。

评价的权威性不能由大模型品牌、引用顶刊数量或双栏模板替代。
本文未建立自然技能的整体好坏标签、人类仲裁一致性或外部基准排名；
不能报告通用质量准确率、语义支持准确率、恶意代码检出率或校准风险上界。
构造故障全检出与自然200包中199包证据不足必须同时保留。
一旦独立验收器遗漏了任务要求，系统仍会出错，故验收器也需要新的反例和外部审查。

\subsection{两个成果各自解决了什么}
成果一给出可使用的判断机制：有机械或行为反证则请求修复，指定输入全部通过才限定放行，无证据或越界则弃权。
它在自建任务中区分了同类高文档分的正确和错误实现；对咨询型自然技能仍不能提供普遍判定。
新增规则与标准库互补，并非所有部分都比现有简单工具更强。

成果二将同一证据用于版本匹配、条件适配和完整路径选择。
相对无质量证据的元数据最短路径，完整产物成功明显增加；
相对逐节点前五截断，AND前提与复合能力保留避免了部分能力缺口。
但已知人工流程和全局风险路径在宽预算下已经完成全部主任务，完整方法没有在这些条件下进一步提高成功数。
版本压力收益也被全局风险路径取得，不能只归于条件分组。
本文没有学到通用路由器，也没有比较已训练的SkillRouter或DSR模型。

本批ECP需要928次探测；只报告执行阶段7.14\%调用节省会误导采用决定。
对已知而稳定的流程，人工预配置在当前批次更省总调用。
网络探索只有在未知配置、多条件复用、维护漂移或紧预算新增覆盖的收益抵消取证成本时才可能值得。
目前没有长期运行、真实token与计费、节点维护时间或端到端延迟数据，不能宣称已实现实际部署降本。

\subsection{有限条件证据无法认证整个任务族}
ECP的平滑频率是排序启发量，不是概率校准成果。
节点错误可相关，上游输出分布会改变下游条件；微小的条件分组无法表达所有输入属性。
主任务成功和单项压力成功都没有阻止组合压力的失败。
这支持在新任务上继续验收和有条件弃权，不支持“开发组无失败即可长期自主执行”。
成果一对未见输入返回越界，成果二则探索新路径；二者的结论强度不同，应在产品中分别显示。

前一阶段18节点为研究者编写，缺陷和任务也是研究者安排，可能使方法有利于已知的条件特征。
四份UCI数据提供公开可追溯输入，但清洗语义、故障和算子跨源共享；两个留出数据集结果完全相同尤其提示模板主导。
这不是对原始1881对象网络的自然端到端改进证明，更不能当作临床或原始分类任务性能。
任务标记由可信适配器给出，自动识别和审计这些标记的成本未测。
小网络穷举没有证明在千级异构网络上的可扩展性。

\subsection{历史实验为何仍保留}
模型引用的高定位有效率只证明来源窗口可追踪，没有建立语义正确性。
300次文本产物的23/240范围内成功，揭示平均逐项通过与完整可用之间的测量差距；
工具受限条件未允许社区脚本运行，故不能推出原技能在正常工具环境下同样失败。
原包加契约和历史自然网络重排没有稳定收益，这些负结果不能被新受控网络成功覆盖或改写。
54及300次请求来自可变服务别名，同一任务重复、多个策略和预算都不是独立样本。
全部条件区间均只对目的性设计作描述，不承担总体显著性或共形风险保证。

\subsection{尚需完成的外部验证}
下一阶段应以独立领域人员的盲标与仲裁检验真实语义质量，公开仲裁前一致性；
将自然技能接入合格隔离执行环境，使用未参与方法设计的真实任务和新的独立验收；
再比较发表实现、同材料模型评审、确定性工具、人工流程及ECP，计入取证和维护的完整成本。
应预先固定来源分组、输入族、预算和重要最小效应，保留拒绝与失败，而不按观察结果选择最有利任务。
此路线是扩大效度的后续任务，不属于本文已达成的结论。
